\documentclass{article}

\usepackage[margin = 1in]{geometry}
\usepackage{amsmath}

\title{Ch. 35 Classwork}
%\author{Principles of Macroeconomics}
%\date{\today}

\begin{document}

%\maketitle
\begin{center}
\textbf{{\Large Ch. 35 Classwork}} 

\bigskip

\end{center}

\begin{enumerate}

\item Suppose the economy is initially in long-run equilibrium. Due to a decline in house prices, suppose that consumers reduce their spending.

	\begin{enumerate}
	
	\item Explain how the decline in consumer spending affects the AD curve.
	
	\item If the central bank wants to stabilize the economy, how will it react to the decline in spending? Use an AD-AS diagram to illustrate and explain your answer.
	
	\item Now, suppose that at the same time that consumer spending falls, oil prices rise sharply due to wars in the Middle East.
	
		\begin{enumerate}
		
		\item Explain how this shock affects the AS curve. How does the immediate short-run equilibrium differ from the short-run equilibrium in part (a)?
		
		\item Discuss, using AD-AS diagrams, the dilemma that the central bank now faces. Keep in mind the Fed's dual-mandate -- keeping inflation low and maintaining full-employment.
		
		\end{enumerate}
	
	\end{enumerate}
	
\item Consider an economy with the following characteristics:

	\begin{align*}
	C &= 2,000 + MPC(Y - T) - 1,000r \\
	I &= 1,000 - 3,000r \\
	G &= 2,200 \\
	NX &= 0 \\
	T &= 1,250 \\
	MPC &= 0.8 \\
	Y^* &= 22,000
	\end{align*}
	
Suppose that the real interest rate, $r$, is initially equal to 5\%, or 0.05.

	\begin{enumerate}
	
	\item Calculate the short-run equilibrium level of real GDP. Is the economy in an expansion or a recession?
	
	\item Suppose that the Fed and the government would like to move the economy back to potential output. The government is only willing to step in with changes in taxes if the Fed is unable to close the gap between potential and actual output. The government is not willing to change spending. Calculate the interest rate and the level of taxes after the necessary actions are taken by the Fed and the government.
	
	\end{enumerate}
	
\item This problem asks you to think about the effects of a sudden contraction in investment and the ensuing consumer pessimism. The twist in this problem is that it captures the realistic feature that taxes are linked to income. The initial position of the economy is as follows.

	\begin{align*}
	C &= 1,000 + MPC(Y - T) \\
	T &= tY \\
	t &= 0.3 \\
	G &= 500 \\
	I &= 1,000 \\
	NX &= 0 \\
	MPC &= 0.9
	\end{align*}

	\begin{enumerate}
	
	\item Calculate the initial level of output in the economy. Assume this also equals potential output.
	
	\item Suppose that weakness in the construction sector of the economy causes planned investment to fall from 1,000 to 900. What is the new level of GDP? Warning -- we can't use the formula for the multiplier that we learned in class because we've changed the model to incorporate income taxes.
	
	\item Further, suppose that the initial decline in investment spending leads to widespread pessimism about the state of the economy, and consumers decide to increase their precautionary saving. Specifically, assume that the marginal propensity to consume falls from 0.9 to 0.8. What is the new level of real GDP?
	
	\item Now suppose that the government wants to use fiscal policy to restore the economy to potential output. By how much would the government need to change the tax rate in order to accomplish this objective?
	
	\item Show each of the changes in parts (b), (c), and (d) on a single, carefully labeled AD-AS diagram.
	
	\end{enumerate}
	
\item In this problem, you will explore the implications of an explicit monetary policy rule that reflects the Fed's dual mandate of pursuing price stability and maximum sustainable employment. To begin, assume that the components of planned aggregate expenditure are given by:

	\begin{align*}
	C &= 3,000 + MPC(Y - T) - 2,500r \\
	MPC &= 0.8 \\
	I &= 1,000 - 2,500r \\
	T &= 1,250 \\
	G &= 1,050 \\
	NX &= 0
	\end{align*}
	
	\begin{enumerate}
	
	\item Compute the equilibrium level of output as a function of the real interest rate. Why should consumption and planned investment depend negatively on the interest rate?
	
	\end{enumerate}
	
	Now suppose that the Fed sets the money supply to achieve an interest rate target. In particular, the Fed sets its interest rate target based on the following rule: \[r = 0.02 + 0.5(\pi - 0.02) \]
	
	\begin{enumerate}
	
	\setcounter{enumii}{1}
	\item Write down an expression showing inflation as a function of current output.
	
	\end{enumerate}
	
	Now suppose that the Fed sets its interest rate target based on the Taylor rule (named after the economist John Taylor). In particular, suppose that: \[r = 0.02 + 0.5(\pi - 0.02) + \frac{1}{2} \left( \frac{Y - Y^*}{Y^*} \right) \]
	
	\begin{enumerate}
	
	\setcounter{enumii}{2}
	\item Assume that the economy is initially at potential output, with inflation equal to its target rate of 0.02. Solve for potential output.
	
	\item Using the value of potential output above, write down an expression showing inflation as a function of current output. Compare this result to the one that you found in part (b).
	
	\item How would the introduction of a Taylor monetary policy rule affect the economy's response to (a) aggregate demand shocks, and (b) aggregate supply shocks? Support your answer with AD-AS diagrams.
	
	\end{enumerate}

\end{enumerate}

\end{document}